% !TEX root = ../algebra_02.tex

\section{Ортогональность в евклидовых пространствах}

\subsection{Процесс ортогонализации Грама-Шмидта}

\begin{Theorem}{Ортогонализация Грама-Шмидта}{}
	Пусть $v_1, \dots, v_n$ — линейно независимая система (ЛИС). Тогда существует ортонормированная система (ОНС) $e_1, \dots, e_n$ такая, что:
	\[
	\mathrm{Lin}(e_1, \dots, e_k) = \mathrm{Lin}(v_1, \dots, v_k) \quad \forall k = 1, \dots, n.
	\]
\end{Theorem}

\begin{proof}
	Доказательство проведем по индукции.
	
	\textbf{База индукции} ($n = 1$): 
	Положим $e_1 = \frac{v_1}{\|v_1\|}$. Очевидно, что $\|e_1\| = 1$ и $\mathrm{Lin}(e_1) = \mathrm{Lin}(v_1)$.
	
	\textbf{Шаг индукции} ($n = k - 1 \to n = k$): 
	Пусть уже построена ОНС $e_1, \dots, e_{k-1}$ для векторов $v_1, \dots, v_{k-1}$. 
	Будем искать новый вектор $e$ в виде:
	\[
	e = v_k + \alpha_1 e_1 + \dots + \alpha_{k-1} e_{k-1}
	\]
	Потребуем, чтобы вектор $e$ был ортогонален всем предыдущим векторам базиса, то есть $(e, e_j) = 0$ для всех $j = 1, \dots, k-1$. Умножим скалярно обе части на $e_j$:
	\[
	(e, e_j) = (v_k, e_j) + \alpha_j (e_j, e_j) = (v_k, e_j) + \alpha_j = 0
	\]
	Отсюда находим коэффициенты: $\alpha_j = -(v_k, e_j)$.
	
	Покажем, что полученный вектор $e \neq 0$. Предположим от противного, что $e = 0$. Тогда:
	\[
	v_k \in \mathrm{Lin}(e_1, \dots, e_{k-1}) = \mathrm{Lin}(v_1, \dots, v_{k-1})
	\]
	Это противоречит условию линейной независимости системы $v_1, \dots, v_n$. Следовательно, $e \neq 0$.
	
	Нормируем вектор $e$, полагая $e_k = \frac{e}{\|e\|}$. Теперь $\|e_k\| = 1$, и $(e_k, e_j) = 0$ для всех $j = 1, \dots, k-1$. 
	Очевидно, что $\mathrm{Lin}(e_1, \dots, e_k) = \mathrm{Lin}(v_k, e_1, \dots, e_{k-1}) = \mathrm{Lin}(v_1, \dots, v_k)$. Шаг индукции доказан.
\end{proof}

\subsection{Ортогональные матрицы}

\begin{Definition}{Ортогональная матрица}{}
	Матрица $U \in M_n(K)$ называется ортогональной, если $U U^T = E_n$, что равносильно $U^T U = E_n$.
	Множество всех ортогональных матриц обозначается как $O_n(K)$.
\end{Definition}

\begin{Lemma}{Свойство множества ортогональных матриц}{}
	Множество $O_n(K)$ образует подгруппу в полной линейной группе $GL_n(K)$.
\end{Lemma}

\begin{proof}
	Пусть $U_1, U_2 \in O_n(K)$. Проверим замкнутость относительно умножения:
	\[
	(U_1 U_2)(U_1 U_2)^T = U_1 U_2 U_2^T U_1^T = U_1 E_n U_1^T = U_1 U_1^T = E_n
	\]
	Если $U \in O_n(K)$, то $U^{-1} = U^T$. Проверим, что обратная матрица также ортогональна:
	\[
	U^T (U^T)^T = U^T U = E_n
	\]
	Следовательно, $O_n(K)$ является подгруппой.
\end{proof}

\begin{Proposition}{Свойства строк и столбцов ортогональной матрицы}{}
	Пусть $U = (a_{ij}) \in M_n(K)$. Следующие условия эквивалентны:
	\begin{itemize}
		\item $U$ — ортогональная матрица.
		\item Строки $U[1, :], \dots, U[n, :]$ образуют ОНС в $M_{1 \times n}(K)$.
		\item Столбцы $U[:, 1], \dots, U[:, n]$ образуют ОНС в $K^n$.
	\end{itemize}
\end{Proposition}

\begin{proof}
	Рассмотрим элемент матрицы $U U^T$ на позиции $(k, l)$:
	\[
	(U U^T)[k, l] = \sum_{j=1}^n U[k, j] \cdot U^T[j, l] = \sum_{j=1}^n a_{kj} \cdot a_{lj} = (U[k, :], U[l, :])
	\]
	Равенство $U U^T = E_n$ выполняется тогда и только тогда, когда скалярное произведение строк равно $1$ при $k=l$ и $0$ при $k \neq l$, то есть когда строки образуют ОНС. Аналогичное рассуждение для $U^T U = E_n$ доказывает утверждение для столбцов.
\end{proof}

\begin{Proposition}{Матрица перехода между ОНБ}{}
	Пусть $E$ — ортонормированный базис (ОНБ) в евклидовом пространстве $V$, а $F$ — некоторый базис $V$. Пусть $U$ — матрица перехода от $E$ к $F$ ($U = M_{E \to F}$). Тогда $F$ является ОНБ тогда и только тогда, когда $U \in O_n(\mathbb{R})$.
\end{Proposition}

\begin{proof}
	Матрица Грама для базиса $F$ выражается через матрицу Грама базиса $E$:
	\[
	\Gamma_F = U^T \Gamma_E U
	\]
	Поскольку $E$ — ОНБ, то $\Gamma_E = E_n$. Следовательно, $\Gamma_F = U^T U$. 
	Базис $F$ является ОНБ $\Leftrightarrow \Gamma_F = E_n \Leftrightarrow U^T U = E_n \Leftrightarrow U \in O_n(\mathbb{R})$.
\end{proof}

\subsection{Ортогональное дополнение}

\begin{Definition}{Ортогональность векторов и ортогональное дополнение}{}
	Два вектора $u, v \in V$ называются ортогональными ($u \perp v$), если их скалярное произведение $(u, v) = 0$.
	
	Пусть $V$ — евклидово пространство, $W \subset V$. Ортогональным дополнением к $W$ называется множество:
	\[
	W^\perp = \{v \in V \mid \forall w \in W: v \perp w\}
	\]
	Очевидно, что $W^\perp$ всегда является линейным подпространством в $V$.
\end{Definition}

\begin{Statement}{Свойства ортогонального дополнения}{}
	Пусть $V$ — евклидово пространство конечной размерности ($\dim V = n < \infty$), а $W, W_1, W_2$ — его подпространства. Тогда:
	\begin{enumerate}
		\item $V^\perp = \{0\}$
		\item $\{0\}^\perp = V$
		\item Если $\dim W = m$, то $\dim W^\perp = n - m$
		\item $V = W \oplus W^\perp$
		\item $W_1 \subset W_2 \Rightarrow W_2^\perp \subset W_1^\perp$
		\item $(W_1 + W_2)^\perp = W_1^\perp \cap W_2^\perp$
		\item $(W_1 \cap W_2)^\perp = W_1^\perp + W_2^\perp$
		\item $(W^\perp)^\perp = W$
	\end{enumerate}
\end{Statement}

\begin{proof}
	Докажем свойства последовательно:
	\begin{itemize}
		\item \textbf{Свойство 1:} Если $v \in V^\perp$, то $v \perp v \Rightarrow (v, v) = 0 \Rightarrow v = 0$.
		\item \textbf{Свойство 2:} Для любого $v \in V$ выполнено $(v, 0) = 0 \cdot (v, 0) = 0$.
		\item \textbf{Свойство 3 и 4:} Пусть $e_1, \dots, e_m$ — ОНБ в $W$. Дополним его до базиса всего пространства векторами $v_{m+1}, \dots, v_n$. Применим процесс ортогонализации Грама-Шмидта и получим ОНБ всего пространства $V$: $(e_1, \dots, e_m, e_{m+1}, \dots, e_n)$.
		
		Для любого $v \in W^\perp$ выполнено $(v, w) = 0$ для всех $w \in W$. В частности, $(v, e_j) = 0$ для $j = 1, \dots, m$. 
		Разложим $v$ по построенному базису: $v = \sum_{i=1}^n \alpha_i e_i$. Из условия ортогональности получаем $\alpha_j = 0$ для $j = 1, \dots, m$. 
		Значит, $v \in \mathrm{Lin}(e_{m+1}, \dots, e_n)$, то есть $W^\perp = \mathrm{Lin}(e_{m+1}, \dots, e_n)$. Отсюда $\dim W^\perp = n - m$.
		
		Если $w \in W \cap W^\perp$, то $(w, w) = 0 \Rightarrow w = 0$. Следовательно, сумма прямая: $W + W^\perp = W \oplus W^\perp$.
		Считаем размерность: $\dim(W \oplus W^\perp) = \dim W + \dim W^\perp = m + (n - m) = n$. Следовательно, $W \oplus W^\perp = V$.
		
		\item \textbf{Свойство 5:} Следует напрямую из определения. Если вектор ортогонален всем элементам б\'ольшего множества $W_2$, то он ортогонален и всем элементам его подмножества $W_1$.
		
		\item \textbf{Свойство 6:} Из свойства 5 имеем $(W_1 + W_2)^\perp \subset W_1^\perp$ и $(W_1 + W_2)^\perp \subset W_2^\perp$, то есть $(W_1 + W_2)^\perp \subset W_1^\perp \cap W_2^\perp$.
		
		Обратно: пусть $v \in W_1^\perp \cap W_2^\perp$. Произвольный вектор $w \in W_1 + W_2$ можно представить как $w = w_1 + w_2$, где $w_1 \in W_1, w_2 \in W_2$. Тогда:
		\[
		(v, w) = (v, w_1) + (v, w_2) = 0 + 0 = 0
		\]
		Значит, $v \in (W_1 + W_2)^\perp$, откуда $W_1^\perp \cap W_2^\perp \subset (W_1 + W_2)^\perp$. Равенство доказано.
		
		\item \textbf{Свойство 8:} Очевидно, что $W \subset (W^\perp)^\perp$, так как любой вектор из $W$ ортогонален любому вектору из $W^\perp$. 
		Посчитаем размерность:
		\[
		\dim(W^\perp)^\perp = n - \dim W^\perp = n - (n - \dim W) = \dim W
		\]
		Так как $W \subset (W^\perp)^\perp$ и их размерности равны, то $W = (W^\perp)^\perp$.
		
		\item \textbf{Свойство 7:} Применим свойство 6 к подпространствам $W_1^\perp$ и $W_2^\perp$, а затем используем свойство 8:
		\[
		(W_1^\perp + W_2^\perp)^\perp = (W_1^\perp)^\perp \cap (W_2^\perp)^\perp = W_1 \cap W_2
		\]
		Взяв ортогональное дополнение от обеих частей, получим $(W_1^\perp + W_2^\perp)^{\perp\perp} = (W_1 \cap W_2)^\perp$, откуда $W_1^\perp + W_2^\perp = (W_1 \cap W_2)^\perp$.
	\end{itemize}
\end{proof}

\subsection{Сопряженные операторы в евклидовом пространстве}

\begin{Definition}{Сопряженный оператор}{}
	Пусть $V$ — евклидово пространство конечной размерности, $A \in \mathrm{End}(V)$ — линейный оператор. Оператор $B \in \mathrm{End}(V)$ называется сопряжённым к $A$, если для любых векторов $v, w \in V$ выполняется равенство:
	\[
	(Av, w) = (v, Bw)
	\]
\end{Definition}

\begin{Proposition}{Существование и единственность сопряженного оператора}{}
	Пусть $A \in \mathrm{End}(V)$. Тогда:
	\begin{enumerate}
		\item Оператор $B$, сопряжённый к $A$, существует и единственен.
		\item Если $E$ — ОНБ в пространстве $V$, то в этом базисе матрица сопряжённого оператора является транспонированной матрицей исходного оператора: если $[A]_E = A$, то $[B]_E = A^T$.
	\end{enumerate}
\end{Proposition}

\begin{proof}
	Зафиксируем произвольный ОНБ $E = (e_1, \dots, e_n)$. Пусть $A = (a_{ij})$ и $B = (b_{ij})$ — матрицы операторов $A$ и $B$ в этом базисе. 
	Любые векторы $v$ и $w$ можно представить через их координаты: $v = E \cdot [v]_E$, $w = E \cdot [w]_E$.
	
	Равенство $(Av, w) = (v, Bw)$ должно выполняться для всех векторов, в частности, для базисных:
	\[
	(Ae_i, e_j) = (e_i, Be_j) \quad \text{для всех } i, j = 1, \dots, n.
	\]
	Вычислим левую часть, зная, что $Ae_i$ — это вектор, координатами которого является $i$-й столбец матрицы $A$: $Ae_i = a_{1i}e_1 + \dots + a_{ni}e_n$.
	\[
	(Ae_i, e_j) = (a_{1i}e_1 + \dots + a_{ni}e_n, e_j) = a_{ji}
	\]
	(так как базис ортонормированный, все скалярные произведения $(e_k, e_j) = 0$ при $k \neq j$, и $(e_j, e_j) = 1$).
	
	Аналогично вычислим правую часть: $Be_j = b_{1j}e_1 + \dots + b_{nj}e_n$, тогда:
	\[
	(e_i, Be_j) = (e_i, b_{1j}e_1 + \dots + b_{nj}e_n) = b_{ij}
	\]
	Следовательно, условие $(Av, w) = (v, Bw)$ равносильно тому, что для всех $i, j$:
	\[
	a_{ji} = b_{ij} \iff B = A^T
	\]
	Это доказывает как единственность сопряжённого оператора (его матрица однозначно определяется как транспонированная), так и его существование (достаточно задать оператор матрицей $A^T$ в ОНБ).
\end{proof}

\subsection{Самосопряженные операторы}

\begin{Statement}{Свойства сопряжённого оператора}{}
	Пусть $A, B \in \mathrm{End}(V)$, а $\lambda \in \mathbb{R}$. Тогда:
	\begin{enumerate}
		\item $(A+B)^* = A^* + B^*$
		\item $(\lambda A)^* = \lambda A^*$
		\item $(AB)^* = B^* A^*$
		\item $(A^*)^* = A$
	\end{enumerate}
	Все эти свойства напрямую следуют из свойств транспонированной матрицы, так как в ОНБ матрица сопряжённого оператора равна транспонированной ($[A^*]_E = [A]_E^T$).
\end{Statement}

\begin{proof}
	Докажем свойство 3. Для любых векторов $v, w \in V$:
	\[
	(v, B^* A^* w) = (Bv, A^* w) = (ABv, w)
	\]
	По определению сопряжённого оператора $(ABv, w) = (v, (AB)^* w)$. Отсюда следует, что $(AB)^* = B^* A^*$.
\end{proof}

\begin{Definition}{Самосопряжённый оператор}{}
	Оператор $A \in \mathrm{End}(V)$ называется самосопряжённым, если $A^* = A$.
	В ортонормированном базисе матрица такого оператора симметрична ($A = A^T$).
\end{Definition}

\begin{Remark}{Связь с билинейными формами}{}
	Если оператор $A$ является самосопряжённым, то отображение $B: V \times V \to \mathbb{R}$, заданное правилом $(v, w) \mapsto (Av, w)$, является симметричной билинейной формой. Действительно:
	\[
	(Aw, v) = (v, Aw) = (v, A^* w) = (Av, w)
	\]
\end{Remark}

\begin{Proposition}{Корни характеристического многочлена}{}
	Пусть $A$ — самосопряжённый оператор. Тогда его характеристический многочлен $\chi_A$ раскладывается на линейные множители над $\mathbb{R}$ (то есть не имеет комплексных корней с ненулевой мнимой частью).
\end{Proposition}

\begin{proof}
	Пусть $\lambda \in \mathbb{C}$ — корень характеристического многочлена $\chi_A$. Зафиксируем ОНБ $E$ и рассмотрим матрицу оператора $A = [A]_E$.
	Рассмотрим оператор умножения на матрицу $A$ в пространстве $\mathbb{C}^n$: $\mathbb{C}^n \xrightarrow{A \cdot} \mathbb{C}^n$. Его характеристический многочлен совпадает с $\chi_A$.
	
	Так как $\chi_A(\lambda) = 0$, существует собственный вектор $x \in \mathbb{C}^n$, $x \neq 0$, такой что $Ax = \lambda x$.
	Рассмотрим следующее произведение и транспонируем его:
	\[
	(Ax)^T \bar{x} = (\lambda x)^T \bar{x} = \lambda x^T \bar{x}
	\]
	С другой стороны, так как матрица $A$ вещественная и симметричная ($A^T = A = \bar{A}$):
	\[
	x^T A^T \bar{x} = x^T A \bar{x} = x^T \overline{A x} = x^T \overline{\lambda x} = \bar{\lambda} x^T \bar{x}
	\]
	Приравнивая два выражения, получаем $\lambda x^T \bar{x} = \bar{\lambda} x^T \bar{x}$.
	Пусть $x = (x_1, \dots, x_n)^T$. Так как $x \neq 0$, имеем:
	\[
	x^T \bar{x} = |x_1|^2 + \dots + |x_n|^2 \neq 0
	\]
	Следовательно, можно сократить на $x^T \bar{x}$, откуда $\lambda = \bar{\lambda} \Rightarrow \lambda \in \mathbb{R}$.
\end{proof}

\begin{Proposition}{Ортогональность собственных векторов}{}
	Собственные векторы самосопряжённого оператора, отвечающие различным собственным значениям, ортогональны.
\end{Proposition}

\begin{proof}
	Пусть $A$ — самосопряжённый оператор, $\lambda, \mu$ — различные собственные значения ($\lambda \neq \mu$), а $v, w$ — соответствующие им собственные векторы.
	Вычислим скалярное произведение $(Av, w)$.
	С одной стороны:
	\[
	(Av, w) = (\lambda v, w) = \lambda (v, w)
	\]
	С другой стороны:
	\[
	(Av, w) = (v, A^* w) = (v, Aw) = (v, \mu w) = \mu (v, w)
	\]
	Приравниваем результаты:
	\[
	\lambda (v, w) = \mu (v, w) \Rightarrow (\lambda - \mu)(v, w) = 0
	\]
	Поскольку собственные значения различны ($\lambda \neq \mu$), необходимо $(v, w) = 0$, то есть $v \perp w$.
\end{proof}

\begin{Theorem}{Спектральная теорема для самосопряжённых операторов}{}
	Пусть $V$ — евклидово пространство, $A \in \mathrm{End}(V)$ — самосопряжённый оператор. Тогда существует ортонормированный базис $E$, в котором матрица оператора $[A]_E$ диагональна.
\end{Theorem}

\begin{proof}
	Доказательство проведем индукцией по размерности пространства $n = \dim V$.
	
	\textbf{База индукции} ($n = 1$): Очевидна.
	
	\textbf{Шаг индукции} ($n - 1 \to n$):
	По доказанному ранее, характеристический многочлен имеет вещественный корень $\lambda$. Следовательно, у оператора $A$ существует собственный вектор $v$, отвечающий этому значению. Нормируем его: $e_1 = \frac{v}{\|v\|}$.
	Рассмотрим ортогональное дополнение $W = \mathrm{Lin}(e_1)^\perp$. Очевидно, что $\dim W = n - 1$.
	Проверим инвариантность подпространства $W$. Пусть $w \in W$, тогда:
	\[
	(Aw, e_1) = (w, A^* e_1) = (w, A e_1) = (w, \lambda e_1) = \lambda (w, e_1) = 0
	\]
	Таким образом, $Aw \perp e_1$, следовательно $Aw \in W$.
	
	Рассмотрим сужение оператора $A_1 = A|_W \in \mathrm{End}(W)$. Очевидно, что оно также является самосопряжённым ($A_1^* = A_1$).
	По индуктивному предположению, в $(n-1)$-мерном пространстве $W$ существует ОНБ $e_2, \dots, e_n$, в котором матрица оператора $A_1$ диагональна.
	Тогда система векторов $e_1, e_2, \dots, e_n$ образует искомый ортонормированный базис всего пространства $V$.
\end{proof}

\subsection{Ортогональные операторы}

\begin{Definition}{Ортогональный оператор}{}
	Оператор $A \in \mathrm{End}(V)$ называется ортогональным, если $A^* A = \mathcal{E}$, где $\mathcal{E}$ — тождественный оператор.
\end{Definition}

\begin{Proposition}{Свойства ортогонального оператора}{}
	Пусть $A \in \mathrm{End}(V)$. Следующие условия эквивалентны:
	\begin{enumerate}
		\item $A$ — ортогональный оператор.
		\item Существует ортонормированный базис $E$, в котором матрица оператора ортогональна: $[A]_E \in O_n$.
		\item В любом ортонормированном базисе матрица оператора ортогональна.
		\item $\forall v, w \in V: (Av, Aw) = (v, w)$ (сохраняет скалярное произведение).
		\item $\forall v \in V: \|Av\| = \|v\|$ (сохраняет длины векторов).
	\end{enumerate}
\end{Proposition}

\begin{proof}
	Докажем эквивалентность последовательно:
	
	\textbf{$1 \Rightarrow 3$}:
	Если $A^* A = \mathcal{E}$, то в любом ОНБ $E$:
	\[
	[A]_E^T [A]_E = [A^*]_E [A]_E = [A^* A]_E = [\mathcal{E}]_E = E_n
	\]
	Следовательно, матрица $[A]_E$ ортогональна.
	
	\textbf{$3 \Rightarrow 2$}:
	Очевидно, так как хотя бы один ОНБ существует по процессу ортогонализации.
	
	\textbf{$2 \Rightarrow 1$}:
	Аналогично $1 \Rightarrow 3$. Из равенства $[A]_E^T [A]_E = E_n$ следует $[A^* A]_E = [\mathcal{E}]_E$, значит $A^* A = \mathcal{E}$.
	
	\textbf{$1 \Rightarrow 4$}:
	\[
	(Av, Aw) = (v, A^* A w) = (v, \mathcal{E} w) = (v, w)
	\]
	
	\textbf{$4 \Rightarrow 1$}:
	По условию для любых $v, w$ выполнено $(Av, Aw) = (v, w)$. Перенесем оператор:
	\[
	(v, A^* A w) = (v, \mathcal{E} w) \Rightarrow (v, (A^* A - \mathcal{E}) w) = 0
	\]
	Так как это выполнено для любого вектора $v$, выберем $v = (A^* A - \mathcal{E}) w$:
	\[
	((A^* A - \mathcal{E}) w, (A^* A - \mathcal{E}) w) = 0
	\]
	Значит, $(A^* A - \mathcal{E}) w = 0$ для всех $w$. Отсюда $A^* A w = w$, то есть $A^* A = \mathcal{E}$.
	
	\textbf{$4 \Rightarrow 5$}:
	Очевидно, достаточно подставить $v = w$.
	
	\textbf{$5 \Rightarrow 4$}:
	Для любых $v, w \in V$ по условию 5 имеем $\|A(v+w)\| = \|v+w\|$. Возведем в квадрат:
	\[
	(A(v+w), A(v+w)) = (v+w, v+w)
	\]
	Раскроем скобки в обеих частях по линейности:
	\[
	(Av, Av) + 2(Av, Aw) + (Aw, Aw) = (v, v) + 2(v, w) + (w, w)
	\]
	Поскольку $(Av, Av) = (v, v)$ и $(Aw, Aw) = (w, w)$, слагаемые сокращаются, и остаётся $(Av, Aw) = (v, w)$.
\end{proof}

\begin{Example}{Матрица поворота}{}
	Рассмотрим матрицу поворота $R_\varphi = \begin{pmatrix} \cos\varphi & \sin\varphi \\ -\sin\varphi & \cos\varphi \end{pmatrix} \in O_2$.
	Её характеристический многочлен:
	\[
	\chi_{R_\varphi}(x) = \begin{vmatrix} \cos\varphi - x & \sin\varphi \\ -\sin\varphi & \cos\varphi - x \end{vmatrix} = (\cos\varphi - x)^2 + \sin^2\varphi = x^2 - 2x\cos\varphi + 1
	\]
	Дискриминант $D/4 = \cos^2\varphi - 1 \le 0$. Следовательно, при $\varphi \neq \pi k$ (где $k \in \mathbb{Z}$) многочлен не имеет вещественных корней.
\end{Example}

\begin{Remark}{}{}
	Можно добавить условие $\varphi \neq \pi k$, чтобы исключить случаи, когда матрица вырождается в скалярную. При $\varphi = 0$ и $\varphi = \pi$ получаем диагональные матрицы:
	\[
	R_0 = \begin{pmatrix} 1 & 0 \\ 0 & 1 \end{pmatrix}, \quad R_\pi = \begin{pmatrix} -1 & 0 \\ 0 & -1 \end{pmatrix}
	\]
\end{Remark}

\begin{Theorem}{Спектральная теорема для ортогональных операторов}{}
	Пусть $A \in \mathrm{End}(V)$ — ортогональный оператор. Тогда существует ортонормированный базис $E$, в котором матрица $[A]_E$ имеет блочно-диагональный вид, причём каждый блок на диагонали имеет вид либо матрицы поворота $R_\varphi$ (при $\varphi \neq \pi k$), либо одномерного блока $(\pm 1)$.
\end{Theorem}